\section{Methodology}
\label{sec:methodology}

\subsection{Interest Profiling and Dense Embeddings}
\label{sec:embeddings}

A user's interest profile is initialized from a descriptive natural language paragraph delineating primary research goals, methodological preferences, and target application domains.

Both publication contents (concatenated title and abstract) and user profiles are projected into a shared semantic vector space using the \texttt{all-MiniLM-L6-v2} sentence transformer:
\begin{equation}
\mathbf{p} = \frac{f_{\text{enc}}(\text{title} \circ \text{" "} \circ \text{abstract})}{\|f_{\text{enc}}(\text{title} \circ \text{" "} \circ \text{abstract})\|_2}, \quad \mathbf{v}_0 = \frac{f_{\text{enc}}(\text{seed\_text})}{\|f_{\text{enc}}(\text{seed\_text})\|_2}
\end{equation}
where $\mathbf{p}, \mathbf{v}_0 \in \mathbb{R}^{384}$ and $\|\mathbf{p}\|_2 = \|\mathbf{v}_0\|_2 = 1.0$. Because all embeddings are strictly $L_2$-normalized upon creation, the semantic cosine similarity reduces to a standard dot product:
\begin{equation}
s(\mathbf{p}, \mathbf{v}) = \cos(\mathbf{p}, \mathbf{v}) = \mathbf{p} \cdot \mathbf{v} = \sum_{j=1}^{384} p_j v_j
\end{equation}

\subsection{Percentile-Gated Relevance Routing}
\label{sec:gating}

Raw cosine similarity scores vary substantially across distinct user interest profiles and scientific subfields; a score of 0.65 may represent an exceptionally close match in a broad domain but a mediocre match in a narrow domain. To achieve robust, domain-invariant decision boundaries, our relevance agent computes thresholds relative to an isolated, fixed \textbf{Reference Sample} $\mathcal{S}_{\text{ref}}$ comprising $M = 5,000$ randomly selected arXiv papers.

Given the active interest vector $\mathbf{v}$, the reference score distribution is computed:
\begin{equation}
\mathcal{D}_{\text{ref}}(\mathbf{v}) = \{ \mathbf{p}_k \cdot \mathbf{v} \mid p_k \in \mathcal{S}_{\text{ref}} \}
\end{equation}
The high and low relevance thresholds, $t_{\text{high}}$ and $t_{\text{low}}$, are calculated as the $P_{\text{high}}$-th and $P_{\text{low}}$-th percentiles of $\mathcal{D}_{\text{ref}}(\mathbf{v})$:
\begin{equation}
t_{\text{high}} = \text{Quantile}\left(\mathcal{D}_{\text{ref}}(\mathbf{v}), \frac{P_{\text{high}}}{100}\right), \quad t_{\text{low}} = \text{Quantile}\left(\mathcal{D}_{\text{ref}}(\mathbf{v}), \frac{P_{\text{low}}}{100}\right)
\end{equation}
The three-way relevance gating rule is defined as:
\begin{equation}
\Phi(\mathbf{p}, \mathbf{v}) = \begin{cases}
\text{HIGH}, & \text{if } s(\mathbf{p}, \mathbf{v}) \ge t_{\text{high}} \\
\text{BORDERLINE}, & \text{if } t_{\text{low}} < s(\mathbf{p}, \mathbf{v}) < t_{\text{high}} \\
\text{LOW}, & \text{if } s(\mathbf{p}, \mathbf{v}) \le t_{\text{low}}
\end{cases}
\end{equation}

\noindent\textbf{Calibration Objective:} Threshold percentiles are calibrated on ground-truth development labels by solving a dual-objective constrained optimization:
\begin{align}
\min \quad & P_{\text{high}} \quad \text{subject to} \quad \text{Precision}_{\text{weighted}}(\text{HIGH}) \ge 0.80 \\
\max \quad & P_{\text{low}} \quad \text{subject to} \quad P_{\text{low}} < P_{\text{high}} \;\land\; \text{MissRate}_{\text{weighted}}(\text{LOW}) \le 0.05 \nonumber
\end{align}
On our calibrated development set, this optimization yields optimal operating levels at $P_{\text{high}} = 98.0$ and $P_{\text{low}} = 90.0$.

\subsection{Asynchronous Human Review Queue and Overflow}
\label{sec:review_queue}

Borderline papers ($t_{\text{low}} < s < t_{\text{high}}$) represent publications with moderate semantic proximity where automated confidence is insufficient. To prevent cognitive fatigue, the review queue enforces strict operational caps:
\begin{itemize}
    \item \textbf{Queue Capacity:} Maximum of 15 borderline tasks per user per monitoring cycle, ordered descending by similarity score.
    \item \textbf{Queue Overflow Policy:} Papers exceeding the 15-task cap are assigned the terminal status \texttt{DEFERRED\_OVERFLOW} and logged to the database without generating review tasks.
    \item \textbf{Task Expiration Policy:} Review tasks remaining undecided after 72 hours (3 calendar days) transition to \texttt{REVIEW\_EXPIRED} and are retired without generating training labels.
    \item \textbf{Idempotent Resume Step:} When a user marks an item as \textbf{Accepted}, a background worker transitions the paper to \texttt{SUMMARIZE\_PENDING}, triggers the summarizer, and appends the paper to the current active digest (\texttt{via='review'}).
\end{itemize}

\subsection{Gated, Interest-Framed LLM Summarization}
\label{sec:summarization}

Unlike standard summarizers that merely condense paper abstracts, our \textbf{Summary Agent} generates \textit{interest-framed} syntheses explaining \textit{why} a specific manuscript matters relative to the researcher's expressed profile.

Summarization is restricted strictly to gate-cleared papers. The agent invokes a hosted LLM API (Google Gemini 1.5 Flash) via a standardized wrapper with low decoding temperature ($T = 0.2$). The prompt structure provides the user's seed interest paragraph and the paper's abstract, instructing the model to generate a structured 3-sentence summary: (1) core methodology, (2) primary empirical result, and (3) concrete relevance to the user's research goals. To prevent duplicate billing, summaries are cached with a unique constraint on \texttt{(user\_id, paper\_id, llm\_model, prompt\_version)}.

\subsection{Multi-Source Feedback and the Effective-Label Rule}
\label{sec:feedback_labels}

The system records two distinct human interaction signals:
\begin{enumerate}
    \item \textbf{Reviewer Decisions:} Captured when a researcher triages borderline abstracts in the review queue ($\text{Accept} \to 1$, $\text{Reject} \to 0$).
    \item \textbf{End-User Reactions:} Captured when a researcher reads a delivered paper summary in the daily research digest ($\text{Like} \to 1$, $\text{Dislike} \to 0$).
\end{enumerate}

Because a user may accept a borderline paper based on its abstract but subsequently dislike the delivered summary, a deterministic conflict resolution rule is required. The database view \texttt{training\_labels} computes one \textbf{effective label} per (user, paper) pair according to the following precedence:
\begin{equation}
y_{\text{eff}}(u, p) = \begin{cases}
y_{\text{user}}(u, p), & \text{if an end-user digest reaction exists} \\
y_{\text{reviewer}}(u, p), & \text{else if a reviewer queue decision exists} \\
\emptyset, & \text{otherwise (skips and expirations yield no label)}
\end{cases}
\end{equation}
Raw events are preserved immutably in an append-only ledger (\texttt{feedback\_events}).

\subsection{Unsupervised Adaptation: Rocchio Centroid Drift}
\label{sec:rocchio}

During the initial cold-start phase ($n < 30$ effective labels), training a supervised parametric classifier is prone to severe overfitting. Instead, the system adapts using a modified, $L_2$-normalized Rocchio centroid formulation~\cite{rocchio1971relevance}:
\begin{equation}
\mathbf{c}_{t+1} = \frac{\alpha \mathbf{v}_0 + \frac{\beta}{|P|} \sum_{\mathbf{p} \in P} \mathbf{p} - \frac{\gamma}{|N|} \sum_{\mathbf{n} \in N} \mathbf{n}}{\left\| \alpha \mathbf{v}_0 + \frac{\beta}{|P|} \sum_{\mathbf{p} \in P} \mathbf{p} - \frac{\gamma}{|N|} \sum_{\mathbf{n} \in N} \mathbf{n} \right\|_2}
\end{equation}
where $\mathbf{v}_0$ is the initial seed vector, $P = \{\mathbf{p} \mid y_{\text{eff}} = 1\}$, $N = \{\mathbf{n} \mid y_{\text{eff}} = 0\}$, and hyperparameters are configured to default values $\alpha = 1.0, \beta = 0.75, \gamma = 0.25$. (If no negative labels exist, the negative summation term is omitted).

The update is triggered upon accumulating 10 effective labels (with $\ge 3$ positives), and retrained after every 25 subsequent feedback events. Whenever a new vector version is generated, the reference sample is immediately re-scored, and updated thresholds ($t_{\text{high}}, t_{\text{low}}$) are computed to preserve fixed percentile shares.

\subsection{Supervised Adaptation: Classifier with Promotion Guard}
\label{sec:supervised_classifier}

Once a researcher accumulates at least 30 effective labels (with $\ge 8$ positive and $\ge 8$ negative examples), the system activates its supervised learning pipeline.

The preference model is formulated as an $L_2$-regularized Logistic Regression classifier~\cite{pedregosa2011scikit}:
\begin{equation}
P(\text{relevant} \mid \mathbf{p}) = \sigma(\mathbf{w}^T \mathbf{p} + b) = \frac{1}{1 + e^{-(\mathbf{w}^T \mathbf{p} + b)}}
\end{equation}
The model is fitted using balanced class weighting to penalize minority-class errors. The regularization hyperparameter $C$ is chosen from $\{0.1, 0.3, 1.0\}$ via internal 5-fold cross-validation.

\noindent\textbf{Autonomous Promotion Guard:} To protect the user from model degradation, a newly trained classifier is evaluated using repeated stratified 5-fold out-of-fold (OOF) cross-validation. A candidate model is promoted to active routing if and only if:
\begin{equation}
\text{ROC-AUC}_{\text{OOF}}(\text{Clf}) \ge 0.65 \;\land\; \text{ROC-AUC}_{\text{OOF}}(\text{Clf}) \ge \text{ROC-AUC}_{\text{OOF}}(\text{Centroid}) + 0.02
\end{equation}
If the candidate fails this promotion guard, the system rejects the model and continues routing via the Rocchio centroid.

Once promoted, routing operates on predicted posterior probabilities:
\begin{equation}
\Phi_{\text{clf}}(\mathbf{p}) = \begin{cases}
\text{HIGH}, & \text{if } P(\text{relevant} \mid \mathbf{p}) \ge p_{\text{high}} \\
\text{BORDERLINE}, & \text{if } p_{\text{low}} < P(\text{relevant} \mid \mathbf{p}) < p_{\text{high}} \\
\text{LOW}, & \text{if } P(\text{relevant} \mid \mathbf{p}) \le p_{\text{low}}
\end{cases}
\end{equation}
with calibrated cutoffs $p_{\text{high}} = 0.70$ and $p_{\text{low}} = 0.30$.
